%%%%%%%%%%%%%%%%%%%%%%%%%%%%%%%%%%%%%%%%%%%%%%%%%%%%%%%%%%%%%%%%%%%%%%%%%%%%%%%%
%%  EINLAGE 3: PROTOKOLL EINES PRAKTIKUMSTAGES   Zwischenabgabe: 09.12.2026 (obligatorisch)
%%  Struktur: Beschreibung (1/4–1/2 S.) + Reflexion (1–2 S.)
%%  Die Laborjournal-Seiten (unredigiert, ALLE des Tages) kommen in den Anhang.
%%%%%%%%%%%%%%%%%%%%%%%%%%%%%%%%%%%%%%%%%%%%%%%%%%%%%%%%%%%%%%%%%%%%%%%%%%%%%%%%
% \Protokolldatum und \Laborbuchlink werden in einstellungen.tex gesetzt

\section{Einlage: Alle Originalseiten ihres Protokolls eines Praktikumstages.}
\label{sec:protokoll}

\begin{hinweise}
Diese Einlage muss aus Ihrem elektronischen Laborjournal sein. Sie verwenden alles was Sie an dem gewählten Tag
aufgeschrieben haben. Sie geben den Link zu Ihrem Laborbucheintrag an.
Verwenden Sie nicht den gleichen Tag wie Ihre Partnerin / Ihr Partner. Machen Sie sich Gedanken über die
folgenden Fragen: Was ist ein gutes Protokoll? Ein gutes Protokoll enthält alle Informationen (Gleichungen,
Literaturhinweise, Handlungsanweisungen, Schaltpläne, Konstruktionszeichnungen, usw.) die notwendig ist, dass
eine andere Person mit Ihrem Bildungsstand, aber ohne weitere Verbindungen zu Ihnen, Ihre Experimente oder
Rechnungen wiederholen kann und Ihre Ergebnisse bewerten kann.
Diese Einlage ist keine Nacherzählung Ihrer Arbeiten an einem Praktikunstag, sondern eine Analyse Ihrer
Dokumentation Ihrer Arbeiten.

\smallskip\emph{Projektbezug:} Praktikumstage sind donnerstags 10--12\,Uhr (O26/233 bzw.\ O26/234).
Abgabe am 09.\,12.\,2026, daher kommen die Tage bis 03.\,12.\ in Frage, z.\,B.\ Platine abholen und
Funktionstest (05.\,11.), Firmware/Kalibrierung des Inertialsensors (12.--26.\,11.) oder die Testmessung
des kompletten Aufbaus (03.\,12.). Tag mit Janosch absprechen (muss verschieden sein), dann
\texttt{\textbackslash Protokolldatum} und \texttt{\textbackslash Laborbuchlink} in \texttt{einstellungen.tex}
setzen.
\end{hinweise}

\subsection{Beschreibung}
\begin{hinweise}
Die Beschreibung gibt alle notwendigen Informationen an, die die Einlage beschreiben, also zum Beispiel das
Datum, die Art des Laborjournals und eine kurze Inhaltsübersicht, aber keine Nacherzählung ($1/4$ Seite, max.
$1/2$ Seite!).
\end{hinweise}
\begin{umfang}{Einlage 3 (Protokoll): Beschreibung}{0.25}{0.5}
Die Einlage umfasst alle unredigierten Seiten meines elektronischen Laborjournals vom \Protokolldatum{}
(Anhang~\ref{anh:protokoll}). Link zum Laborbucheintrag: \expandafter\url\expandafter{\Laborbuchlink}.

\platzhalter{Art/Software des elektronischen Laborjournals (gemeinsam mit \Partner, Arbeitspaket 4.1),
kurze Inhaltsübersicht des Tages (keine Nacherzählung)}
\end{umfang}

\subsection{Reflexion}
\begin{hinweise}
In der Reflexion sollen Sie sich Gedanken machen über die Wirkung dieser Einlage und der Arbeit mit ihr auf
Ihren Lernfortschritt.
\begin{itemize}\itemsep0pt
\item In welchem Zusammenhang steht die Einlage mit Ihrem Schwerpunktthema und Ihrem Projektthema?
\item Verstehen Sie noch Ihren Notizen?
\item Was haben Sie bei der Erstellung dieser Einlage gelernt?
\item Wie hat Ihnen die Einlage beim Verständnis Ihres Projektthemas geholfen?
\item Wo gab es Schwierigkeiten bei der Beschäftigung mit der Einlage?
\item Was wünschen sie sich von Ihren Notizen beim erneuten Lesen?
\item Was hätten Sie anders gemacht, wenn Sie die Einlage nochmals erstellen müssten?
\item Was sagen Ihre Kommilitoninnen oder Kommilitonen dazu, als Kommentar oder wenn Sie mit Ihnen über die Einlage diskutieren?
\item Haben Sie in Ihrem Protokoll logische oder physikalische Fehler entdeckt?
\item Gab es bei der Beschäftigung mit der Einlage eine überraschende neue Erkenntnis?
\end{itemize}
\textbf{Die Reflexion für eine Einlage sollte mindestens eine Seite und maximal 2 Seiten lang sein. Die
Längenangaben sind ohne Titel, Bilder, Tabellen und Tabellen- oder Bildlegenden zu verstehen.}
Tipp: Konkrete Stellen im Journal benennen (\enquote{Seite 2 oben fehlt die Verstärkereinstellung \dots}).
\end{hinweise}
\begin{umfang}{Einlage 3 (Protokoll): Reflexion}{1}{2}
\platzhalter{Analyse der eigenen Dokumentation: Was ist ein gutes Protokoll? Wäre mein Experiment
wiederholbar? Leitfragen erkennbar beantworten}
\end{umfang}

%\subsection{Eigene Bewertungskriterien}
%...
